\documentclass[10pt]{article}

\usepackage[utf8]{inputenc}

\usepackage[T1]{fontenc}

\usepackage{amsmath}

\usepackage{amsfonts}

\usepackage{amssymb}

\usepackage[version=4]{mhchem}

\usepackage{stmaryrd}

\usepackage{graphicx}

\usepackage[export]{adjustbox}

\graphicspath{ {./images/} }

\usepackage{mathrsfs}



\begin{document}

«сколь угодно малые" числа $\varepsilon>0$ заменяются «сколь угодно мелкими» покрытиями $\omega$. Это позволяет аналогичную теорему сформулировать для любых нормальных пространств и придти к заключению, что в некром (все же наглядном) геометрическом смысле всякое $n$-мерное нормальное пространство «похоже» и даже «сколь угодно мало отличается» от $n$-мерного полиэдра.



Одной из важнейших теорем Р. т. является т. н. те орема о с ущ е с твен ных о тоб р а жен и я х, лежащая в основе значительной части всей этой теории. Пусть $f$ - непрерывное отображение (нормального) пространства $X$ на $n$-мерный шар $Q^{n}$ с границей $S^{n-1}$. Пусть ФСX-прообраз сферы $S^{n-1}$ при этом отображении, $\Phi=f^{-1} S^{n-1}$. Отображение $f: X \rightarrow Q^{n}$ наз. с у щ е т в е н н ы м, если всякое непрерывное отображение $g: X \rightarrow Q^{n}$, совпадающее с $f$ во всех точках $x \in \Phi$, есть отображение на весь пар $Q^{n}$. Упомянутая т е о р е м а А ле к с а н д р о в а утверждает, что нормальное пространство $X$ тогда и только тогда имеет размерность $\operatorname{dim} X \geqslant n$, когда пространство $X$ можно существенно отобразить на $n$-мерный пар. Из этой теоремы выводится теорема суммы (доказанная для компактов П. С. Урысоном и К. Менгером еще в самом начале развития $\mathrm{P}$. т.): если (нормальное) пространство $X$ размерности $\operatorname{dim} X=n$ является объединением конечного или счетного числа своих замкнутых множеств $\Phi_{k}$, то по крайней мере для одного из этих $\Phi_{k}$ имеет место $\operatorname{dim} \Phi_{k}=n$.



Теорема о существенных отображениях лежит в основе т. н. Гомоло гической те о ри и р а зм е р н о с т и, позволяющей применить к изучению размерности в весьма общих предположениях методы алгебраич. топологии. Понятие гомологической размерности связано с понятиями цикла и гомологии и поэтому предполагает, что наряду с топологич. пространством $X$ дана и нек-рая коммутативная группа (\$), к-рая наз. группой коэффициентов. Тогда можно говорить о циклах компакта $X$ по этой группе коэфффи-



\includegraphics[max width=\textwidth, center]{2023_07_08_a6a2494a29728a6cb415g-1(2)}

циклах, гомологичных нулю в $X$ по области коәффициентов Њ, причем эти понятия можно әквивалентным образом понимать как в смысле теории гомологий Александрова - Чеха, так и в смысле гомологич. теории Вьеториса.



После әтого можно определить гомологич. размерность компакта $X$ по группе коэфффициентов (\$) как наибольпее целое число $n$ такое, что в компакте $X$ имеется $(n-1)$-мерный цикл $Z^{n-1}$, гомологичный нулю в $X$, но не гомологичный нулю на нек-ром своем носителе Ф. Оказывается, что $\operatorname{dim} X$ есть гомологич. размерность по группе $x$, являющейся факторгруппой группы всех действительных чисел по подгруппе целых чисел, и она является напбольшей среди всех вообще гомологич. размерностей.



Если от циклов и гомологий перейти к коциклам и когомологиям, то получится когомологическая размерность, причем когомологич. размерность компакта по данной дискретной группе $\mathfrak{A}$ равна гомологич. размерности по бикомпактной группе $\mathfrak{B}$, двойственной группе $\mathfrak{U}$ в смысле теории характеров Понтрягина. Отсюда следует, что размерность $\operatorname{dim} X$ совпадает с когомологической размерностью по группе целых чисел.



$\begin{aligned} & \text { Лит. см. при ст. Размерность. } \text { П. С. Александров. } \\ & \text { РАЗМЕРНОСТИ ФУНКЦИЯ - целочисленная функ- }\end{aligned}$



\includegraphics[max width=\textwidth, center]{2023_07_08_a6a2494a29728a6cb415g-1(3)}

удовлетворяющая условиям: 1) $d(x+y)+d(x y)=d(x)+$ $+d(y)$ для любых $x, y \in L ; 2)$ если $[x, y]$-простой интервал в $L$, то $d(y)=d(x)+1$. Для решетки, все ограниченные цепи к-рой конечны, существование Р. Ф. эквивалентно дедекиндовости әтой решетки. Существует и болеө общеө определение Р. Ф. на ортомодулярной рөшөтке или ортомодулярном тастично упорядоченном множестве, где значениями Р. ф. могут быть произвольные действительные числа или даже функции (см. [3]).



Лит.: [1] С к о р н я к о в Л. А., Элементы теории структур, M., 1970; [2] B i r k h of $\mathbf{f}$ G., Lattice theory, 3 ed., Providence, 1967; [3] K a 1 m b a c h G., Omolattices, L., 1981 .



PABMEPHOCT Т. Фобанова РАЗМЕРНОСТНЫИ ИНВАРИАНТ - целое число $d(X)$, определяемое для каждого топологич. пространства $X$ данного класса $\mathscr{K}$ и обладающее достаточным количеством свойств, сближающих его с обычным числом измерений евклидовых многомерных пространств. При этом от класса $\mathscr{K}$ требуется, чтобы он содержал все кубы любого числа измерений и вместе с любым данным пространством $X$, являющимся его элементом, содержал в качестве элемента и всякое пространство, гомеоморфное пространству $X$. От Р. и. $d(X)$ предполагается, во всяком случае, что для гомеоморфных пространств $X$ и $X^{\prime}$ всегда $d(X)=d\left(X^{\prime}\right)$ и что для $n$-мерного куба $I^{n}$ выполняется $d\left(I^{n}\right)=n$. Среди Р. и. важнейшими являются т. н. к л а с с и че с ки е ра зм е р н о сти-Лебега размерность $\operatorname{dim} X$ и (большая и малая) индуктивная размерность Ind $X$, ind $X$.



Имеют место следующие предложения, выделяющие размерность $\operatorname{dim} X$ среди всех Р. и., определенных соответственно в классе всех (метрических) компактов, всех метризуемых и всех сепарабельных метризуемых пространств и тем решающие для этих пространств проблему аксиоматич. определения размерности. Единственный, удовлетворяющий нижеперечисленным условиям 1), 2), 3), определенный в классе $\mathscr{K}$ всех (метрических) компактов $X$ размерностный инвариант $d(X)$ есть размерность $\operatorname{dim} X=$ Ind $X=$ ind $X$ (т е о р е м a А ле к с ан д р о в а).



Условие 1) (аксио ма П у анка ре). Если пространство принадлежит к классу $\mathscr{K}$ п $d(X)$ равно неотрицательному целому числу $n$, то в $X$ содержится замкнутое подпространство $X_{0}$, для к-рого $d\left(x_{0}\right)<n$ и такое, что множество $X \backslash X_{0}$ несвязно.



Условие 2) (акси им а конечно й с умм ы ). Если принадлежащее к классу $\mathscr{K}$ пространство $X$ есть объединение двух замкнутых подпространств $X_{1}$ и $X_{2}$, для к-рых $d\left(X_{1}\right) \leqslant n, d\left(X_{2}\right) \leqslant n$, то и $d(X) \leqslant n$. У с ловие 3) (аксиома Б ра у э р а в ме т р и ч е с к о й фо рм е). Если $X$ - принадлежащее к классу $\mathscr{\kappa}$ пространство и $d(X)$ есть неотрицательное целое число $n$, то существует такое положительное число $\varepsilon$, что для всякого пространства $Y$, являющегося образом пространства $X$ при каком-либо



\includegraphics[max width=\textwidth, center]{2023_07_08_a6a2494a29728a6cb415g-1(1)}

этом отображение $f$ компакта $X$ на компакт $Y$ наз. є-о т о б р а ж е ни е м, если оно непрерывно, и полный прообраз $f^{-1}(y)$ каждой точки $y \in Y$ имеет в $X$



\begin{center}

\includegraphics[max width=\textwidth]{2023_07_08_a6a2494a29728a6cb415g-1}

\end{center}



Т е о рем а IЩ е п и н а [2]. Размерность $\operatorname{dim} X$ есть единственный, определенный соответственно в классе $\mathscr{K}$ всех метрических, всех сепарабельных метрич. пространств $X$ размерностный инвариант $d(X)$, удовлетворяющий следующим условиям (т е о р е м а IЩ е п и на).



Условие 1) (аксиома П у анка ре, см. выше).



Условие 2) (аксиома счетной сумм ы). Если принадлежащее к классу $\mathscr{K}$ пространстро $X$ есть объединение счетного числа своих замкнутых подпространств $X_{k}, k=1,2, \ldots$, для каждого из к-рых $d\left(X_{k}\right) \leqslant n$, то и $d(X) \leqslant n$.



Условие 3) (акси има Б р а у э р а в о бШе й ф о р ме). Если для принадлежащего к классу ¿̋ прос'гранства $X$ имеется $d(X) \leqslant n$, то существует та-





\end{document}